%! Characteristics of Rational Functions — Unit 4, Lesson 4.0 — Group Activity (Answer Key)
\documentclass[10pt]{article}
\usepackage{algebra2-article}
\usepackage{algebra2-key}   % pulls in -boxes; do NOT also load -boxes
\newcommand{\TallMath}[1]{$\displaystyle #1\rule[-1.4em]{0pt}{3.2em}$}
\newcommand{\lgrid}[4]{%
  \draw[step=1, linegray!40, very thin] (#1,#3) grid (#2,#4);
  \draw[->, semithick, charcoal] (#1,0)--(#2,0) node[below right, font=\scriptsize]{$x$};
  \draw[->, semithick, charcoal] (0,#3)--(0,#4) node[left, font=\scriptsize]{$y$};}

\begin{document}

\pageheader{Unit 4, Lesson 4.0}{Group Activity --- Answer Key}
\namepartnerperiod

\vspace{0.08in}

\begin{headlinebox}{forestbg}
\small\textbf{Reading Rational Graphs.} Every characteristic of a rational function is something you
can \emph{point to} on its graph --- or predict from its denominator. With your partner, read each
graph carefully and \emph{justify} your answers. Write every asymptote as an \textbf{equation}
($x=3$, $y=0$), never as a bare number.
\end{headlinebox}

\vspace{0.06in}

\begin{tcolorbox}[colback=white, colframe=black!40, title=\textbf{Tier R --- Remediate},
    fonttitle=\bfseries, arc=1.5mm, left=3mm, right=3mm, top=2mm, bottom=2mm, breakable]
The graph shows $g(x) = \dfrac{2}{x+2}$.
\begin{center}
\begin{tikzpicture}[scale=0.46]
  \lgrid{-6}{4}{-5}{5}
  \draw[navy, dashed, semithick] (-2,-5)--(-2,5);
  \node[navy, font=\scriptsize, above] at (-2,5) {$x=-2$};
  \begin{scope}
    \clip (-6,-5) rectangle (4,5);
    \draw[very thick, forest, domain=-6:-2.4, samples=80, smooth] plot (\x,{2/((\x)+2)});
    \draw[very thick, forest, domain=-1.6:4, samples=80, smooth] plot (\x,{2/((\x)+2)});
  \end{scope}
  \fill[charcoal] (0,1) circle (3pt) node[above right, font=\scriptsize]{$(0,1)$};
\end{tikzpicture}
\end{center}
\begin{enumerate}[leftmargin=1.7em, itemsep=7pt]
  \item The graph comes in \ans{2} pieces. The input that is missing is $x=$ \ans{-2},
        because it makes the \ans{denominator} equal to zero.
  \item \textbf{Vertical asymptote:} \ans{$x=-2$} \quad \textbf{Horizontal asymptote:}
        \ans{$y=0$} \quad (write each as an \emph{equation})
  \item \textbf{Domain:} \ans{$(-\infty,-2)\cup(-2,\infty)$ --- all reals except $x=-2$}
  \item \textbf{$y$-intercept:} \ans{$(0,1)$}. \quad Does $g$ have an \textbf{$x$-intercept}?
        \ans{No.} \; Explain, using the numerator. \ans{The numerator is $2$, never $0$.}
  \item On each piece the curve falls, so $g$ is \textbf{decreasing} on \ans{$(-\infty,-2)$} and on
        \ans{$(-2,\infty)$}. \quad $g(x)>0$ on \ans{$(-2,\infty)$}.
\end{enumerate}
\end{tcolorbox}

\begin{tcolorbox}[colback=white, colframe=black!40, title=\textbf{Tier A --- Approaching Proficiency},
    fonttitle=\bfseries, arc=1.5mm, left=3mm, right=3mm, top=2mm, bottom=2mm, breakable]
Two rational functions --- \emph{but only one of them has asymptotes}. Fill in the whole table, then
answer the justification question.
\begin{center}
\begin{tikzpicture}[scale=0.42]
  % p(x) = (x-2)/(x+1)
  \begin{scope}
    \lgrid{-6}{5}{-5}{6}
    \draw[navy, dashed, semithick] (-1,-5)--(-1,6);
    \draw[navy, dashed, semithick] (-6,1)--(5,1);
    \begin{scope}
      \clip (-6,-5) rectangle (5,6);
      \draw[very thick, forest, domain=-6:-1.6, samples=90, smooth]
        plot (\x,{((\x)-2)/((\x)+1)});
      \draw[very thick, forest, domain=-0.5:5, samples=90, smooth]
        plot (\x,{((\x)-2)/((\x)+1)});
    \end{scope}
    \fill[charcoal] (2,0) circle (3pt); \fill[charcoal] (0,-2) circle (3pt);
    \node[charcoal, font=\scriptsize] at (-0.5,-6.4){$p(x)=\dfrac{x-2}{x+1}$};
  \end{scope}
  % q(x) = (x^2-1)/(x-1)
  \begin{scope}[xshift=13.5cm]
    \lgrid{-4}{4}{-3}{5}
    \begin{scope}
      \clip (-4,-3) rectangle (4,5);
      \draw[very thick, forest] (-4,-3)--(4,5);
    \end{scope}
    \draw[forest, very thick, fill=white] (1,2) circle (4.2pt);
    \node[charcoal, font=\scriptsize] at (0,-4.4){$q(x)=\dfrac{x^2-1}{x-1}$};
  \end{scope}
\end{tikzpicture}
\end{center}
\vspace{0.16in}
\renewcommand{\arraystretch}{1.5}
\begin{tabularx}{\linewidth}{@{}>{\bfseries}p{0.30\linewidth} X X@{}}
 & \textbf{$p(x)=\dfrac{x-2}{x+1}$} & \textbf{$q(x)=\dfrac{x^2-1}{x-1}$} \\
\hline
Excluded value(s) & \ans{$x=-1$} & \ans{$x=1$} \\
Vertical asymptote & \ans{$x=-1$} & \ans{none} \\
Horizontal asymptote & \ans{$y=1$} & \ans{none ($n>m$)} \\
Hole (if any) & \ans{none} & \ans{$(1,2)$} \\
$x$-intercept & \ans{$(2,0)$} & \ans{$(-1,0)$} \\
$y$-intercept & \ans{$(0,-2)$} & \ans{$(0,1)$} \\
Increasing or decreasing? & \ans{increasing on both branches} & \ans{increasing} \\
\end{tabularx}
\vspace{4pt}

\textbf{Justify.} Factor $q(x)=\dfrac{x^2-1}{x-1}$ completely and simplify it. Then explain why
$x=1$ produces a \emph{hole} rather than a \emph{vertical asymptote} --- and why $x=1$ is
\emph{still} excluded from the domain even after you cancel.
\ansline{$q(x)=\frac{(x-1)(x+1)}{x-1}=x+1$ for $x\neq 1$. The factor $(x-1)$ \emph{cancels}, so the}
\ansline{denominator no longer blows up near $x=1$: nearby outputs close in on $2$, not on $\pm\infty$.}
\ansline{But the \emph{original} denominator is $0$ at $x=1$, and cancelling is only legal when $x\neq1$, so $1$ can never re-enter the domain --- the graph is the line $y=x+1$ with the single point $(1,2)$ punched out.}
\end{tcolorbox}

\begin{tcolorbox}[colback=white, colframe=black!40, title=\textbf{Tier E --- Extension},
    fonttitle=\bfseries, arc=1.5mm, left=3mm, right=3mm, top=2mm, bottom=2mm, breakable]
\textbf{Part 1 --- No graph allowed.} For each function, factor where you can, then fill in the row.
Write ``none'' when there is nothing to report.
\begin{center}
\renewcommand{\arraystretch}{1.6}
\begin{tabularx}{\linewidth}{@{}l X X X X@{}}
\hline
\textbf{Function} & \textbf{Excluded value(s)} & \textbf{Vert.\ asymptote(s)} & \textbf{Hole} & \textbf{Horiz.\ asymptote} \\
\hline
\TallMath{\frac{3}{x-5}} & \ans{$x=5$} & \ans{$x=5$} & \ans{none} & \ans{$y=0$} \\
\TallMath{\frac{2x}{x+4}} & \ans{$x=-4$} & \ans{$x=-4$} & \ans{none} & \ans{$y=2$} \\
\TallMath{\frac{x^2-25}{x-5}} & \ans{$x=5$} & \ans{none} & \ans{$(5,10)$} & \ans{none} \\
\TallMath{\frac{x+1}{x^2+3x+2}} & \ans{$x=-1,\,-2$} & \ans{$x=-2$} & \ans{$(-1,1)$} & \ans{$y=0$} \\
\hline
\end{tabularx}
\end{center}
Exactly one of the four has \emph{no} horizontal asymptote. Which one, and what is it about the
degrees that causes that?
\ansline{$\frac{x^2-25}{x-5}$. The numerator's degree ($2$) is \emph{greater} than the denominator's ($1$) ---}
\ansline{top-heavy --- so the outputs grow without bound instead of settling toward a level. (Its graph is the line $y=x+5$ with a hole at $(5,10)$.)}

\vspace{5pt}
\textbf{Part 2 --- What an asymptote \emph{means}.} A club pays a one-time \$250 screen-printing setup
fee plus \$6 per shirt. The \textbf{average cost per shirt} for $n$ shirts is
\[
  A(n) = \frac{6n + 250}{n}.
\]
\begin{center}
\renewcommand{\arraystretch}{1.25}
\begin{tabular}{|c|c|c|c|c|c|}
\hline
$n$ (shirts) & $10$ & $25$ & $50$ & $125$ & $250$ \\ \hline
$A(n)$ (dollars) & $31$ & $16$ & $11$ & $8$ & $7$ \\ \hline
\end{tabular}
\end{center}
\begin{enumerate}[label=(\alph*), leftmargin=1.9em, itemsep=6pt]
  \item Use the degree comparison to find the horizontal asymptote of $A$. \ans{$y=6$} \\
        What does that number \emph{mean} for the club?
        \ansline{The degrees tie ($n=m=1$), so the level is the ratio of leading coefficients, $6/1=6$.}
        \ansline{The more shirts the club orders, the closer the average cost per shirt gets to \$6 --- the true printing cost of one shirt.}
  \item The average cost gets closer and closer to that level but never equals it. Explain why, in
        terms of the \$250 setup fee.
        \ansline{$A(n)=6+\frac{250}{n}$: the setup fee is always split among the shirts, so there is always}
        \ansline{an extra $250/n$ dollars on top of \$6. That term shrinks toward $0$ but is never $0$, so the average stays just above \$6.}
  \item $n=0$ is an excluded value, so $A$ has a vertical asymptote at $n=0$. Explain why $n=0$ would
        be meaningless here \emph{even if} the algebra allowed it, and state a sensible domain for
        this model.
        \ansline{You cannot compute a cost \emph{per shirt} when no shirts are made --- there is nothing to}
        \ansline{average. A sensible domain is the positive integers, $n\ge 1$ (a whole number of shirts).}
\end{enumerate}
\end{tcolorbox}

\begin{teachernote}
\textbf{Tier placement.} Tier~R is pure graph-reading on a one-asymptote graph; put students there
who wrote bare numbers instead of equations in the notes. Tier~A is the discriminating task of the
lesson --- $q$ \emph{looks} like a line, and groups who fill in ``vertical asymptote: $x=1$'' for it
have not yet internalized the cancel test. Tier~E is equation-only work plus the A2.F.2f
interpretation beat.

\textbf{What to listen for.} In Tier~A, the phrase you want is ``the factor cancels, so it is a hole''
--- not ``there is no asymptote because it is a line.'' Push groups to say \emph{why} the line has a
gap. In Tier~E Part~2(b), the algebra $A(n)=6+\frac{250}{n}$ is the whole explanation; if a group is
stuck, have them do that division and the horizontal asymptote becomes obvious.

\textbf{Common errors.} Tier~R item~4: students often invent an $x$-intercept from the graph's
flattening arm; the numerator argument is the fix. Tier~E row~3: expect ``VA at $x=5$'' out of habit
--- send them back to factor first. Tier~E row~4: expect only \emph{one} excluded value; both $-1$ and
$-2$ come from the original denominator, and only one of them survives as an asymptote.
\end{teachernote}

\end{document}
